\documentclass[10pt,a4paper]{amsart}
\usepackage{amsmath,amssymb,amsthm,graphicx}
% For commutative diagrams we use Xy-pic
\usepackage[all]{xy}

% We use 2cell for 2-commutative diagrams.
\xyoption{2cell}
% Theorem environments.
%
\theoremstyle{plain}
\newtheorem{theorem}[subsection]{Theorem}
\newtheorem{proposition}[subsection]{Proposition}
\newtheorem{lemma}[subsection]{Lemma}

\theoremstyle{definition}
\newtheorem{definition}[subsection]{Definition}
\newtheorem{example}[subsection]{Example}
\newtheorem{exercise}[subsection]{Exercise}
\newtheorem{situation}[subsection]{Situation}

\theoremstyle{remark}
\newtheorem{remark}[subsection]{Remark}
\newtheorem{remarks}[subsection]{Remarks}

\numberwithin{equation}{subsection}

% Macros
%
\def\lim{\mathop{\mathrm{lim}}\nolimits}
\def\colim{\mathop{\mathrm{colim}}\nolimits}
\def\Spec{\mathop{\mathrm{Spec}}}
\def\Hom{\mathop{\mathrm{Hom}}\nolimits}
\def\Ext{\mathop{\mathrm{Ext}}\nolimits}
\def\SheafHom{\mathop{\mathcal{H}\!\mathit{om}}\nolimits}
\def\SheafExt{\mathop{\mathcal{E}\!\mathit{xt}}\nolimits}
\def\Sch{\mathit{Sch}}
\def\Mor{\mathop{\mathrm{Mor}}\nolimits}
\def\Ob{\mathop{\mathrm{Ob}}\nolimits}
\def\Sh{\mathop{\mathit{Sh}}\nolimits}
\def\NL{\mathop{N\!L}\nolimits}
\def\CH{\mathop{\mathrm{CH}}\nolimits}
\def\proetale{{pro\text{-}\acute{e}tale}}
\def\etale{{\acute{e}tale}}
\def\QCoh{\mathit{QCoh}}
\def\Ker{\mathop{\mathrm{Ker}}}
\def\Im{\mathop{\mathrm{Im}}}
\def\Coker{\mathop{\mathrm{Coker}}}
\def\Coim{\mathop{\mathrm{Coim}}}

% Boxtimes
%
\DeclareMathSymbol{\boxtimes}{\mathbin}{AMSa}{"02}

%
% Macros for moduli stacks/spaces
%
\def\QCohstack{\mathcal{QC}\!\mathit{oh}}
\def\Cohstack{\mathcal{C}\!\mathit{oh}}
\def\Spacesstack{\mathcal{S}\!\mathit{paces}}
\def\Quotfunctor{\mathrm{Quot}}
\def\Hilbfunctor{\mathrm{Hilb}}
\def\Curvesstack{\mathcal{C}\!\mathit{urves}}
\def\Polarizedstack{\mathcal{P}\!\mathit{olarized}}
\def\Complexesstack{\mathcal{C}\!\mathit{omplexes}}
% \Pic is the operator that assigns to X its picard group, usage \Pic(X)
% \Picardstack_{X/B} denotes the Picard stack of X over B
% \Picardfunctor_{X/B} denotes the Picard functor of X over B
\def\Pic{\mathop{\mathrm{Pic}}\nolimits}
\def\Picardstack{\mathcal{P}\!\mathit{ic}}
\def\Picardfunctor{\mathrm{Pic}}
\def\Deformationcategory{\mathcal{D}\!\mathit{ef}}
\usepackage{cleveref}
\crefname{originallemma}{original lemma}{original lemmas}
\Crefname{originallemma}{Original lemma}{Original lemmas}
\makeatletter
\newlabel{lemma-localize-NL}{{00S7}{source}{}{sourcepointer-1}{}}
\newlabel{lemma-localize-NL@cref}{{[originallemma][0][]00S7}{[1][0][]source}{}{}{}}
\makeatother
\crefname{originallemma}{original lemma}{original lemmas}
\Crefname{originallemma}{Original lemma}{Original lemmas}
\makeatletter
\newlabel{lemma-NL-localize-bottom}{{07BS}{source}{}{sourcepointer-2}{}}
\newlabel{lemma-NL-localize-bottom@cref}{{[originallemma][0][]07BS}{[1][0][]source}{}{}{}}
\makeatother
\title{Smoothness of a ring map is preserved by localizing the source and by inverting an element of the base}
\author{The Stacks Project Authors}
\date{stacks-project:algebra.tex:local-snapshot-20261012; GNU FDL 1.2 or later}
\begin{document}
\maketitle
\noindent\textit{Editorial note.} All delivered mathematical text is the author's own, copied from the pinned Stacks Project chapter file algebra.tex; no statement or proof step was rewritten, reordered or removed.\par\medskip
\noindent\textit{Editorial note.} This is the first card this lane authors from stacks-algebra, the source newly assigned to it, and the first card in the L10000-L11999 range.\par\medskip
\noindent\textit{Editorial note.} The staged candidate list distributed with this source records ext\_refs = 0 for this label, but the proof does reference two same-chapter lemmas; both are recorded as navigation stubs whose tags were looked up by the same-chapter suffix rule, and the cotangent-complex argument stays inside the counted span.\par\medskip
\noindent\textit{Editorial note.} Assembly history: r001 was rejected with "assembly differs from reviewed statement locators" because the hand-written body\_spans carried shortened labels while the contract (prepared\_delivery.py line 249) requires body\_spans to equal the top-level statement/proof spans entry by entry; r002 rebuilds body\_spans from those spans so the two are identical. No span coordinate, statement or proof text changed.\par\medskip
\noindent\textit{Allowed standard prerequisites (editorial).}\begin{itemize}
\item The naive cotangent complex of a ring map. Same chapter, Lemma labelled lemma-localize-NL; invoked at proof line 37339. Identifies the naive cotangent complex of S\_g over R with the base change of that of S over R.
\item Localisation of the naive cotangent complex in the base. Same chapter, Lemma labelled lemma-NL-localize-bottom; invoked at proof line 37345. Gives the analogous statement for the second assertion, after inverting an element of the base.
\end{itemize}\medskip
\ifdefined\hypertarget\hypertarget{sourcepointer-1}{}\fi
\noindent\textit{Source reference.} Original lemma 00S7: sources/stacks-algebra/pinned-original-source.tex, source lines 35909--35909. Navigation pointer for \textbackslash{}ref\{lemma-localize-NL\} used inside the counted proof (line 35909); the cotangent-complex argument that consumes it is retained in the span.\par\medskip
\ifdefined\hypertarget\hypertarget{sourcepointer-2}{}\fi
\noindent\textit{Source reference.} Original lemma 07BS: sources/stacks-algebra/pinned-original-source.tex, source lines 35807--35807. Navigation pointer for \textbackslash{}ref\{lemma-NL-localize-bottom\} used inside the counted proof (line 35807); the cotangent-complex argument that consumes it is retained in the span.\par\medskip
\begin{lemma}
\label{lemma-localize-smooth}
Let $R \to S$ be a smooth ring map.
Any localization $S_g$ is smooth over $R$.
If $f \in R$ maps to an invertible element of $S$,
then $R_f \to S$ is smooth.
\end{lemma}

\begin{proof}
By Lemma \ref{lemma-localize-NL} the naive cotangent
complex for $S_g$ over $R$ is the base change of the naive cotangent
complex of $S$ over $R$. The assumption is that the naive cotangent
complex of $S/R$ is $\Omega_{S/R}$ and that this is a finite projective
$S$-module. Hence so is its base change. Thus $S_g$ is smooth over $R$.

\medskip\noindent
The second assertion follows in the same way from
Lemma \ref{lemma-NL-localize-bottom}.
\end{proof}

\end{document}
